% SPDX-License-Identifier: MIT
\chapter{Listen, Tabellen und Bäume mit Model/View}\label{chap:models}
\ChapterSummary{Sie erfahren, wie Daten und ihre Darstellung getrennt
werden, wie Selektion und sichtbarer Ausschnitt funktionieren und
weshalb ein ListView nicht für jede Datenzeile ein eigenes Widget
erzeugen muss. Die Beispiele beschreiben den derzeitigen M5-Core-
Vertrag und seine Grenzen.}

\section{Warum nicht tausend Labels?}
Eine Administrationsoberfläche kann hunderte Hosts, Dienste,
Firewallregeln oder Logzeilen anzeigen müssen.
Würde jede Zelle zu einem eigenen Widget, würden
Lebensdauer, Messung, Selektion und Redraw unnötig teuer.
Das Toolkit setzt deshalb für strukturierte Daten auf
\term{Model/View}: Das Modell kennt Inhalt und Änderungen,
die View kennt Interaktion und einen sichtbaren Ausschnitt,
und ein SelectionModel verwaltet die aktuelle Auswahl.

\begin{center}
\begin{tikzpicture}[node distance=13mm,every node/.style={draw,
 rounded corners=2pt,minimum width=34mm,minimum height=11mm,
 fill=SASDPale,font=\sffamily\small,align=center}]
\node (model) {Datenmodell\\Inhalt + Revision};
\node[right=of model] (selection) {SelectionModel\\gewählter Eintrag};
\node[below=of selection,xshift=-22mm] (view) {View\\Viewport + Ereignisse};
\draw[->,thick,SASDBlue] (model)--(selection);
\draw[->,thick,SASDBlue] (model)--(view);
\draw[->,thick,SASDBlue] (selection)--(view);
\end{tikzpicture}
\end{center}

\section{ListModel und ListView}
Ein \term{ListModel} bietet \texttt{rowCount()}
und \texttt{textAt(row)}. Das mitgelieferte
\term{StringListModel} speichert Strings und
kennt \texttt{append}, \texttt{insert}, \texttt{erase},
\texttt{setText} und \texttt{reset}.
Bei inhaltlichen Änderungen werden
Modellbeobachter informiert; die Revisionsnummer
steigt. Das \term{ListSelectionModel} verwaltet
eine optionale Zeilenauswahl und kann vorwärts
oder rückwärts navigieren.

\begin{lstlisting}[language=SASDCpp,caption={Sichtbaren Listenausschnitt an ein Modell binden}]
StringListModel items{{"dns", "firewall", "logging"}};
ListSelectionModel selected;
selected.setModel(&items);

ListView view;
view.setModel(&items);
view.setSelectionModel(&selected);
view.setViewport(0, 10);

selected.select(1);
for (const auto& row : view.visibleRows()) {
    // row.row, row.text, row.selected auswerten
}
\end{lstlisting}
Dieses Beispiel ist reine Core-Datenlogik;
es erzeugt ohne Presentation-Sink noch
keine sichtbare Terminal-Anwendung.
In der realen Oberfläche gehört die View
in den Containerbaum; Modell und
SelectionModel leben ausreichend lange.

\section{Ein Viewport ist ein Vertrag, kein Datenverlust}
Mit \texttt{setViewport(...)} bestimmen Sie den
voraussichtlich sichtbaren Zeilenausschnitt.
\texttt{visibleRows()} liefert dafür
Werte-Snapshots statt direkter Besitzrechte
an interne Modell-Strings.
Das Modell selbst bleibt vollständig erhalten.
Nach einer Modellmutation wird der Viewport
normalisiert; eine Auswahl soll in einem
sinnvollen sichtbaren Bereich bleiben.

\section{TableModel: Rechteckige Daten}
Ein \term{TableModel} stellt \texttt{rowCount},
\texttt{columnCount}, \texttt{headerAt(column)}
und \texttt{textAt(row,column)} bereit.
\term{StringTableModel} verwaltet eine
rechteckige Tabelle: Jede Zeile muss zur
Spaltenzahl passen. Die View kann
sowohl Zeilen als auch Spalten virtualisieren;
\term{TableSelectionModel} wählt eine optionale
\term{TableCell} aus Zeile und Spalte.

\begin{lstlisting}[language=SASDCpp]
StringTableModel table(
    {"Host", "Status"},
    {{"fw01", "ok"}, {"dns01", "offline"}});
TableSelectionModel selection;
selection.setModel(&table);
TableView view;
view.setModel(&table);
view.setSelectionModel(&selection);
view.setViewport(0, 12, 0, 2);
selection.select(0, 1);

const auto headers = view.visibleHeaders();
const auto rows = view.visibleRows();
\end{lstlisting}
Das Beispiel kann unverändert in Core-Tests
oder im Datenaufbau einer realen Form
verwendet werden. Eine produktive
Darstellung von Tabellenkanten,
Spaltenbreiten oder Scrollbars bleibt eine
Aufgabe der jeweiligen Presentationsschicht.
Ein \term{TableView} ist nicht automatisch
ein kompletter spreadsheetartiger Editor.

\section{TreeModel: Hierarchie mit Pfaden}
Ein \term{TreeNodePath} enthält eine Folge
von Kindindizes. Das ist ein strukturierter
Pfad, \emph{keine} dauerhaft unveränderliche
Datenbank-ID. Ein \term{TreeModel} beantwortet
\texttt{childCount(parent)} und
\texttt{textAt(node)}. Ein
\term{StringTreeModel} kann Wurzeln,
Kindknoten, Umbenennen und Reset verwalten.
Die \term{TreeView} kann Knoten
expandieren, einklappen und nur den
sichtbaren Teilbaum zurückgeben.

\begin{lstlisting}[language=SASDCpp]
StringTreeModel tree{{"Netz", "System"}};
const TreeNodePath root{{0}};
const TreeNodePath child = tree.append(root, "DNS");

TreeSelectionModel choice;
choice.setModel(&tree);
TreeView view;
view.setModel(&tree);
view.setSelectionModel(&choice);
view.expand(root);
view.setViewport(std::nullopt, 15);
choice.select(child);
\end{lstlisting}
Bei strukturellen Änderungen darf die
Anwendung einen alten Pfad nicht
unbesehen als weiterhin identisch ansehen.
Für real persistierte Objekte benötigt
die Domänenschicht stabile Kennungen,
die sie mit Pfaden abgleichen kann.

\section{Beobachtungen, Besitz und Lebensdauer}
Modelle sind \term{Component}-Objekte,
die Views halten aber nur lebensdauersicher
\emph{nichtbesitzende} \term{Reference}-Handles.
Die \term{Subscription}-Tokens melden sich
bei der Zerstörung automatisch ab.
Das verhindert hängende rohe Zeiger
in Beobachterlisten, ersetzt aber
keine sorgfältige Anwendungstopologie.
Benutzeraktionen, die nach Ablauf eines
Models eintreffen, müssen kontrolliert
ins Leere laufen.

\section{Aktueller Implementationsumfang}
Die Core-Verträge für List-, Table- und
Tree-Modelle, SelectionModels und Views
sind implementiert. Terminal- und
Rendered-Präsentationen sind je nach
Widget und Demo-Fälligkeit unterschiedlich
weit integriert. Planen Sie deswegen vor
einem Produktivprojekt einen sichtbaren
End-to-End-Test des konkret benötigten
Widget/Backend-Paars ein. Aus der Existenz
eines Core-Headers folgt noch nicht,
dass jede Oberflächenvariante die gesamte
Bedienung schon anbietet.

\PracticalTip{Für SASD Gatehold wäre
eine Firewall-Regelliste ein gutes
TableModel; der markierte Regel-Datensatz
wird über ein SelectionModel identifiziert.
Zustand und Berechtigung zum Anwenden
gehören in die Domänenlogik und die
gebundenen Commands, nicht in den
Zeilenrenderer.}
